\chapter{Conclusiones}

Al concluir el ciclo de desarrollo y validación técnica de la plataforma SaaS orientada a la digitalización de las PYMES gastronómicas de Portoviejo, se logró comprobar que la aplicación estratégica de una solución integrada (gestión operativa e inteligencia artificial) es técnicamente viable y usable en entorno controlado. Este proyecto evidencia que es posible desarrollar una plataforma propia adaptada al contexto local que potencialmente puede reducir la dependencia de herramientas dispersas y de plataformas externas con comisiones elevadas. La validación técnica realizada sienta las bases para futuras investigaciones que evalúen la transformación de procesos manuales y fragmentados en flujos digitales centralizados en operación real. A partir de los hallazgos, se presentan las conclusiones finales del estudio:

Los resultados obtenidos durante el desarrollo y la validación técnica de la plataforma en entorno controlado permiten afirmar que la integración de los módulos de pedidos (con pagos en efectivo y transferencia), inventarios, vista de cocina en móvil (aceptar o rechazar pedidos, notificar cuando esté listo para recoger), administración de repartidores y dashboard analítico alcanzó un nivel de funcionalidad adecuado durante las pruebas con usuarios de la PYME piloto. La centralización de la información en tiempo real en las tareas guiadas demostró el potencial de la plataforma para reducir tiempos de registro de pedidos y errores operativos, en coherencia con los umbrales de éxito definidos en la metodología. Asimismo, el enfoque modular (backend NestJS, aplicación móvil React Native + Expo con vistas cliente, repartidor y cocina, sitio web Next.js para administración y compra como cliente) permitió un despliegue estable en entorno de staging (VPS, Vercel) y una experiencia de uso coherente para usuarios de prueba. Estos datos evidencian que una plataforma SaaS integrada constituye una solución técnicamente viable y usable, sentando las bases para futuras investigaciones que evalúen su adopción e impacto en operación real frente a los procesos manuales predominantes en las PYMES gastronómicas de Portoviejo.

Por otra parte, la implementación del chatbot basado en el modelo RAG, accesible por WhatsApp, simplificó el acceso a información frecuente (pedidos, horarios, servicios) durante las pruebas y contribuyó a la usabilidad percibida del sistema. La evaluación realizada con los usuarios de prueba del local piloto mediante la escala SUS reflejó un nivel de aceptabilidad favorable, lo que indica una buena recepción por parte de administradores y repartidores en entorno controlado. Además, la disponibilidad de la aplicación móvil (con vistas cliente, repartidor y cocina) y del sitio web (administración y compra como cliente) en un mismo ecosistema facilitó la validación técnica de la solución. Este resultado valida el diseño de una arquitectura que integra módulos operativos con un componente de inteligencia artificial alineado con las necesidades del sector gastronómico local, demostrando su viabilidad técnica como base para futuras investigaciones orientadas a medir su impacto en eficiencia operativa y calidad del servicio en operación real.

Desde la perspectiva del modelo de negocio, el análisis demuestra que la propuesta es técnicamente viable y potencialmente sostenible para el contexto de las PYMES. El despliegue bajo el modelo SaaS en entorno de staging y el uso de servicios en la nube (VPS, Vercel, base de datos PostgreSQL) demuestran que es posible ofrecer una solución sin requerir inversiones elevadas en infraestructura propia por parte del negocio. Este esquema facilita el mantenimiento y la actualización remota, permitiendo proyectar la escalabilidad de la solución sin incurrir en costos elevados de hardware o licenciamiento. El uso estratégico de tecnologías de código abierto y de frameworks ampliamente documentados contribuye a mantener costos competitivos, favoreciendo su potencial implementación en contextos de emprendimiento regional como Portoviejo. Futuras investigaciones deberán evaluar la viabilidad económica y la sostenibilidad financiera en operación real a mayor escala.

Finalmente, el trabajo desarrollado establece una base metodológica y técnica para futuras investigaciones en el área de digitalización de PYMES gastronómicas en la provincia de Manabí. Los resultados de la validación técnica en entorno controlado confirman que la combinación de una plataforma SaaS modular con un componente de inteligencia artificial (chatbot RAG por WhatsApp) es técnicamente factible y usable en soluciones orientadas a la gestión operativa y a la atención al cliente. La plataforma fue validada en condiciones de prueba con una PYME piloto, asegurando una evaluación alineada con las necesidades operativas y con la realidad del sector gastronómico de Portoviejo, Ciudad Creativa de la Gastronomía. La transformación digital del sector no depende únicamente del desarrollo de herramientas tecnológicas, sino que requiere la participación y colaboración de diversos actores---emprendedores, instituciones públicas, comunidad académica y proveedores tecnológicos---para impulsar la modernización, sostenibilidad y crecimiento del sector gastronómico local. Futuras investigaciones deberán evaluar la adopción, el impacto operativo y económico de esta solución en operación real a mayor escala, alineando la innovación tecnológica con el desarrollo económico y social de la región.
